\begin{table}[htbp] \centering \caption{Summary Statistics Indonesia}\label{tab:summary_stats_IDN_unb_nonag} \begin{threeparttable} \begin{tabular}{l cccc} \toprule
 &  &  &  & $t$-test \\
\midrule
Non-Agricultural &  &   33.9\% &   66.1\% &  \\  \addlinespace \addlinespace
Log Consumption & 12.05 & 11.75 & 12.20 & -0.46*** \\
 & (0.79) & (0.74) & (0.78) &  \\  \addlinespace
Log Income & 14.89 & 14.41 & 15.09 & -0.68*** \\
 & (1.14) & (1.15) & (1.08) &  \\  \addlinespace
Female & 0.44 & 0.43 & 0.45 & -0.02*** \\
 & (0.50) & (0.49) & (0.50) &  \\  \addlinespace
Age (years) & 39.94 & 43.40 & 38.16 & 5.24*** \\
 & (13.61) & (15.02) & (12.45) &  \\  \addlinespace
Education (years) & 7.97 & 5.54 & 9.22 & -3.67*** \\
 & (4.63) & (3.92) & (4.47) &  \\  \addlinespace
Household Size & 4.65 & 4.55 & 4.70 & -0.15*** \\
 & (2.07) & (1.98) & (2.12) &  \\  \addlinespace


\cmidrule{2-5}
Observations & 93,038 & 46,797 & 69,289 &  \\  \addlinespace
Individuals & 29,716 &  &  &  \\
Non-switchers &    6.6\% &  &  &  \\  \addlinespace
\bottomrule \end{tabular} \begin{tablenotes}[flushleft] \footnotesize \item{Summary statistics for Indonesia for the unbalanced panel across all five waves. Source: IFLS. The table reports means and standard deviations (in parentheses) based on individual-year pairs. See section 3 for further details. All variables have the same number of observations, except for income, which is missing for some observations. Income has 61,300 observations.} \end{tablenotes} \end{threeparttable} \end{table}